\addchap{Preface}

This book began with a list of textbooks that should exist but do not. One entry read, in effect: biology and physics students take a course in experimental method before their experiments matter, but deep learning researchers learn to run experiments by word of mouth. Nobody teaches them how much the variance between seeds has to be before a difference means something, how an unequal tuning budget can turn a comparison around, how to fit a scaling law and how far to extrapolate it, or how to design an ablation that does not confound. The entry proposed to teach all of this through the field's own famous mistakes.

The book that resulted has fifteen chapters in five parts, on noise, fair comparison, causes, scale and measurement. Each chapter opens with a conclusion that the field once believed and that turned out to be wrong, or much narrower than claimed: that adaptive optimizers generalize worse than plain gradient descent, that large batches generalize worse because they find sharp minima, that batch normalization works by reducing covariate shift, that model size should grow much faster than data, that some abilities of language models emerge abruptly with scale. The chapter reconstructs the experiment behind the conclusion, reports what later work found, and names the methodological failure. It then teaches the method that would have caught it, with the statistics and mathematics that make the method precise, reproduces the failure and its fix in a small experiment that runs on a laptop in minutes, and ends with a protocol for the next experiment of the same kind.

The intended readers are graduate students and researchers who train networks, in machine learning and in the many sciences that now use it, and anyone who reads experimental papers in deep learning and wants to know how much to believe them. The prerequisites are undergraduate probability and a first course in statistics, linear algebra and calculus, and some practice in training networks. Each chapter builds the statistics it uses, and an appendix collects the tools that several chapters share. The prologue explains the plan of the book.

The experiments in this book are small on purpose. A network trained on eight-by-eight digits cannot show everything that a language model trained for a month can; but every failure examined here is a failure of design, not of scale, and a design can be learned on a small experiment and carried to a large one. The reader who runs the scripts, and writes down a prediction before each one, will learn more from them than from the text.

\paragraph{A note on authorship.} This book was written by Claude (Opus 5.5), an AI model made by Anthropic, and conceived and edited by Xiaoyu Li, who proposed the subject and its plan. Every claim attributed to a paper was checked against the paper or a reliable account of it wherever these could be reached, every number from an experiment of this book was produced by a script that the reader can run, and the colophon at the end describes how the book was made. Mistakes surely remain. Readers who find one are invited to treat it as the book treats the field's: as a precise statement of what went wrong, and a pointer to what comes next.
